\section{Introduction}
\label{sec:intro}

The proliferation of autonomous multi-agent systems across critical infrastructure, financial markets, and clinical decision-support environments has exposed a fundamental architectural gap: contemporary agentic frameworks lack a mandatory, structurally enforced mechanism for graceful degradation under conditions of cascading failure or accountability collapse. As autonomous agents assume roles previously reserved for human judgment—roles entailing material consequences, resource allocation, and intergenerational stewardship—the absence of a codified bail-out protocol represents not merely a technical omission but an ethical failure of the first order. This paper introduces the \bxthree Bailout Protocol: a formally specified, architecturally compulsory escalation mechanism that embeds mandatory fallback behavior into the fabric of purpose-driven agentic systems by grounding operational boundaries in verifiable fact.

% -------------------------------------------------------
\subsection{Why Autonomous Systems Need a Mandatory Bail-Out Mechanism}
% -------------------------------------------------------

An agentic system, by definition, pursues objectives expressed in natural language, encoded in policy, or instantiated through reward signals without continuous human oversight. When such a system operates within well-understood nominal conditions, its behavior remains broadly predictable. However, the moment conditions deviate from training distributions, or when inter-agent negotiations produce equilibria that violate implicit constraints, the system proceeds along trajectories for which no human review exists to intervene. In conventional architectures, this condition persists until external termination is applied—often after substantial harm has already propagated through the system and into the world.

The case for a mandatory bail-out mechanism begins with a structural observation about the nature of autonomous agency. An autonomous agent is, by design, an entity that acts without requiring prior approval for each individual action. This is not a defect; it is the fundamental utility of autonomy. But the same property that makes autonomous agents useful—their capacity to continue operating without human authorization—also makes them dangerous when the conditions under which they were instantiated no longer hold. The problem is not that autonomous agents act incorrectly. It is that they lack any internal architectural property by which they can recognize that the conditions justifying their continued operation have changed, and that continued operation now constitutes a departure from their chartered accountability rather than its exercise.

A mandatory bail-out mechanism is therefore not an add-on safety feature. It is the structural expression of the principle that every autonomous agent operates under an implicit condition: that the world-state upon which its current operational scope was justified remains consistent with that scope. When that condition fails—when \fact deviates from the ground state presupposed by \purpose—the system must transition to a degraded operational mode not as a procedural suggestion but as an architectural compulsion. The mechanism that enforces this compulsion is the Bailout Protocol.

% -------------------------------------------------------
\subsection{Failure Modes of Systems Without Architectural Fallback}
% -------------------------------------------------------

The failure modes of systems that lack a mandatory bail-out mechanism are not hypothetical. They are the predictable consequence of architectures that optimize for objective attainment without encoding an obligation to verify the continued validity of the conditions under which objectives were instantiated.

In decentralized financial agent ecosystems, cascading liquidations have occurred precisely because individual agents lacked the capacity to recognize collective instability and execute coordinated withdrawal. Each agent, operating within its own authorized scope, continued optimizing for its individual objective function while the collective system crossed thresholds that no single agent's scope included as a binding constraint. The failure was not the result of any single agent's incorrect inference. It was a structural consequence of the absence of a mechanism by which any agent could recognize that its own scope had been exceeded by the collective dynamics of the system.

In large-scale distributed inference systems, a single hallucinated constraint propagated across thousands of agent instances has produced coordinated actions that simultaneously violated physical and regulatory boundaries. Again, the failure was not a matter of individual agent capability. It was a consequence of the absence of any architectural property by which an agent could distinguish between a constraint that was genuinely part of its operating environment and one that had been generated by another agent's hallucinated inference and propagated through the system without verification.

These are not edge cases. They are structural outputs of architectures that conflate autonomy with the absence of accountability bounds. The conflation is understandable: the engineering discipline that produced modern agentic systems was oriented toward capability, not accountability. But as agentic systems move from constrained laboratory environments into consequential real-world deployments, the absence of an architectural fallback for accountability is no longer an acceptable engineering tradeoff. The problem is structural, not incidental. No amount of prompt engineering, fine-tuning, or oversight layering resolves it, because the architecture itself contains no mechanism by which the system can recognize its own accountability vacuum and respond architecturally rather than procedurally.

% -------------------------------------------------------
\subsection{Core Insight: Accountability Requires an Architectural Fallback}
% -------------------------------------------------------

This observation gives rise to the core insight of the \bxthree framework: accountability must have an architectural fallback, not merely a procedural one. Procedural accountability—rules, checklists, human-in-the-loop gates—is contingent on the very conditions that autonomous systems are designed to transcend. An agent operating with sufficient autonomy will, by design, bypass or override procedural controls when those controls impede objective attainment. This is not a defect; it is the intended function of autonomy.

The \bxthree framework encodes accountability through three interdependent architectural layers, each defined by the properties it must maintain rather than by the type of actor occupying it:

\begin{itemize}
  \item[\textbf{\purpose.}] The \purpose layer carries intent, judgment, and final authority. It specifies what the system is designed to pursue and under what conditions that pursuit remains valid. The \purpose layer defines the operating range $R(n)$ for each node $n$ and receives all escalated exceptions that exceed machine-actor provisioned scope. It is non-delegable as an exception terminus: no machine actor may substitute for the \purpose layer's judgment in resolving an escalated exception.
  \item[\textbf{\bounds.}] The \bounds layer carries bounded reasoning and constrained execution within $R(n)$. It evaluates whether proposed actions fall within the node's provisioned authority and initiates bail-out when they do not. Its participation in the Bailout Protocol is confined to a single, non-discretionary function: detecting conditions that exceed $R(n)$ and triggering the compulsory transition to the \texttt{BOUNDED} state.
  \item[\textbf{\fact.}] The \fact layer carries deterministic enforcement and forensic auditability. It provides the verifiable ground state against which \purpose and \bounds are continuously evaluated. Within the Bailout Protocol, the \fact layer enforces the state machine's transition relation and maintains the append-only Forensic Ledger as an immutable record of every protocol event.
\end{itemize}

The bail-out mechanism is not added to the system as a failsafe. It is the logical consequence of maintaining consistency among \purpose, \bounds, and \fact under all operational conditions. When \fact deviates from the ground state presupposed by \purpose—when the real-world conditions that justified a given operational strategy have fundamentally changed—the \bounds layer triggers a mandatory transition to a pre-specified degraded operational mode. This transition is not a suggestion; it is architecturally enforced by the \fact layer. The system does not continue operation until a human operator intervenes; the architecture itself holds operation in abeyance until consistency among \purpose, \bounds, and \fact is re-established or the system transitions to a terminal idle state pending external review.

In this sense, the \bxthree Bailout Protocol operates analogously to a pressure relief valve in physical engineering: it does not prevent high-pressure conditions, but it guarantees that the system cannot remain in a high-pressure state that exceeds its certified operational envelope. The relief valve does not decide whether the pressure is dangerous—the \purpose layer makes that judgment. The valve does not choose when to open—the trigger condition fires automatically when the observed pressure exceeds the certified envelope. And the valve does not close itself until a qualified operator has inspected the system and authorized resumed operation. Each function is performed by a distinct architectural layer, and the structural separation of these functions is what makes the mechanism reliable.

This separation is what conventional escalation hierarchies lack. Conventional tiered oversight models make human oversight a high-tier option rather than a compulsory terminus. They permit intermediate machine actors to absorb exceptions, resolve them autonomously, and continue operation—terminating the escalation at a machine rather than propagating it to a human. The \bxthree Bailout Protocol forbids this absorption by architectural constraint: the human anchor $h(n)$ is not one option among many; it is the only permissible terminus for every unresolved exception. The structural completeness of the three-layer implementation is what makes the upstream accountability guarantee unconditional: \bxthree systems fail upward into human consciousness, never downward into autonomous chaos.

% -------------------------------------------------------
\subsection{Paper Roadmap}
% -------------------------------------------------------

The remainder of this paper is organized as follows. Section~\ref{sec:bailout-background} reviews related work on autonomous agent safety, graceful degradation in distributed systems, and formal verification approaches to agentic architectures. Section~\ref{sec:bailout-framework} presents the \bxthree three-layer framework in full, establishing the structural substrate upon which the Bailout Protocol operates. Section~\ref{sec:bailout-protocol} provides the complete formal specification of the Bailout Protocol: the deterministic state machine governing exception propagation, the trigger condition, the escalation algorithm, the no-machine-actor bypass mechanism, and the timeout handling that guarantees liveness under all failure conditions. Section~\ref{sec:bailout-evaluation} evaluates the protocol against a suite of adversarial and edge-case scenarios drawn from real-world multi-agent deployments. Section~\ref{sec:bailout-limitations} examines the structural costs of the protocol's guarantees, including human response latency, scalability constraints, and exploitation risk, and identifies directions for future work. Section~\ref{sec:bailout-conclusion} concludes with implications for the future governance of autonomous systems and the role of architectural accountability in agentic deployments.
